\documentclass[10pt]{article}
\usepackage[margin=0.85in]{geometry}
\usepackage{amsmath,amssymb,amsthm}
\newtheorem{theorem}{Theorem}
\title{Equal consecutive Ulm quotients: a counterexample to 00000004273}
\author{earthking11}
\date{October 6, 2026}
\begin{document}
\maketitle
\section*{Statement and actual invariants}
The conjecture describes the Ulm cardinal sequence of a countable
$p$-group and requires realizable sequences to decrease strictly in
ordinal order. We exhibit a finite abelian $2$-group with two equal,
nonzero consecutive Ulm invariants. Neither terminal zeroes nor an
infinite-stage convention is involved.

Take the actual additive group
$G=\mathbb Z/2\mathbb Z\oplus\mathbb Z/4\mathbb Z$.
It is finite, hence countable, and $4x=0$ for every $x\in G$; therefore it
is $2$-primary. Write
$U_n=(2^nG)[2]$, the subgroup of elements annihilated by $2$ inside the
image of multiplication by $2^n$. At each finite stage the corresponding
Ulm quotient is $Q_n=U_n/U_{n+1}$. We use its actual cardinality; the
usual vector-space-dimension convention gives the same equal adjacent
values in this example.

\begin{theorem}
The actual Ulm quotient cardinalities of this countable $2$-group do
not form a strictly decreasing sequence: $|Q_0|=|Q_1|=2$.
\end{theorem}
\begin{proof}
In the indicated coordinate groups,
\[
\begin{split}
 U_0&=\{(0,0),(1,0),(0,2),(1,2)\},\\
 U_1&=\{(0,0),(0,2)\},\qquad U_2=\{(0,0)\}.
\end{split}
\]
The first-coordinate map $U_0\to\mathbb Z/2\mathbb Z$ is a surjective
homomorphism with kernel $U_1$. The map
$U_1\to\mathbb Z/2\mathbb Z$ sending $(0,0)$ to $0$ and $(0,2)$ to $1$
is a surjective homomorphism with kernel $U_2$. The actual quotient
groups $Q_0$ and $Q_1$ are therefore both isomorphic to
$\mathbb Z/2\mathbb Z$, by the first isomorphism theorem.
Their cardinals are equal and nonzero. Strict anti-monotonicity would
require $|Q_1|<|Q_0|$, a contradiction.
If Ulm entries are instead recorded as dimensions over $\mathbb F_2$,
these same quotients are both one-dimensional, so strict decrease still
fails. The formal capstone uses literal cardinal-valued entries.
\end{proof}

An ordinal-indexed strictly decreasing sequence has strictly decreasing
restriction to its indices $0$ and $1$. Failure on this finite pair thus
refutes the proposed necessary condition for realizability, regardless
of how later ordinal stages or a ``full'' sequence are presented. The
additional lattice-closure clause is not needed.

\section*{Formal verification and earlier submission}
Lean uses Mathlib's actual additive group \texttt{ZMod 2 $\times$ ZMod 4},
homomorphism images and kernels, additive subgroups, and actual quotient
groups. It verifies the subgroup descriptions, the two surjective
coordinate homomorphisms, their kernels, quotient equivalences, and
both natural and literal \texttt{Cardinal.mk} cardinalities.
The final theorem simultaneously proves countability, the $2$-primary
condition, and failure of standard \texttt{StrictAnti} for the actual
cardinal sequence.

Earlier closed PR \#309 used this same finite counterexample, but its
capstone encoded the supposed strict decrease as a tailored prohibition
on two quotient certificates rather than directly negating strict
decrease of a cardinal-valued sequence. This replacement supplies that
semantic bridge, uses Mathlib's quotient objects, and credits the prior
example; it does not claim that PR \#309 received a rejecting review.
The development compiles on Lean 4.33.1 with pinned Mathlib. It has no
unfinished proofs, native evaluation, or extra axioms. The only printed
dependencies are propositional extensionality, classical choice, and
quotient soundness. The tiny coordinate checks range over eight group
elements, not a long brute-force computation.
\end{document}
